\chapter{L'environnement de recherche}\label{ch:pl:the-research-environment}

En 2013, trois économistes qui répliquaient un article influent sur la dette publique et la croissance ont obtenu le tableur de travail des auteurs. Une erreur de codage
dans celui-ci excluait cinq pays --- l'Australie, l'Autriche, la Belgique, le Canada et le Danemark --- de l'analyse~; une fois cette erreur et deux autres
problèmes corrigés, la croissance moyenne des pays dont la dette publique dépassait 90~\% du PIB s'établissait à 2.2~\%, et non à $-0.1~\%$ comme publié (Herndon,
Ash et Pollin). L'économie n'était pas ce qui avait fait défaut. Ce qui avait fait défaut, c'était un environnement dans lequel un calcul ne pouvait être rejoué depuis le début
que par son auteur, et même pas par lui une fois que l'état du tableur s'était écarté de ce qu'il affichait. Les
notebooks d'une plateforme de recherche ont la même propriété, et plus encore. Ce chapitre audite un notebook qui ment de la même façon, le rejoue depuis le début,
élève son calcul au rang de bibliothèque, et reproduit le chiffre corrigé à partir d'un répertoire propre.

\section{Notebooks et état caché}

\begin{definition}[Notebook de calcul]\label{def:pl:the-research-environment:notebook}
Un \emph{notebook de calcul}\index{notebook de calcul} est un document composé de cellules --- de code, de texte et des sorties produites par le code ---
exécutées une par une dans un interpréteur vivant (le noyau), dans l'ordre que l'utilisateur choisit~; le format de Jupyter le stocke en JSON, avec, pour chaque
cellule de code, sa source, ses sorties enregistrées et le compteur d'exécution du noyau au moment de sa dernière exécution (un entier, ou null si elle n'a jamais été exécutée).
\end{definition}

\begin{definition}[État caché]\label{def:pl:the-research-environment:hidden}
Un notebook présente un \emph{état caché}\index{état caché} lorsque ses sorties enregistrées dépendent de quelque chose que ses cellules, exécutées depuis le début, ne
reproduiraient pas~: une cellule exécutée hors d'ordre, une cellule exécutée deux fois, une cellule modifiée ou supprimée après son exécution, un nom défini seulement à l'invite de commande.
\end{definition}

Le notebook est le moyen le plus rapide d'explorer des données, et le plus facile pour publier un chiffre que personne ne peut recalculer. Pimentel, Murta, Braganholo et
Freire ont collecté des notebooks Jupyter sur GitHub et ont tenté de les exécuter~: sur 863\,878 exécutions tentées de notebooks valides, seules 24.11~\%
se sont déroulées sans erreur et seulement 4.03~\% ont produit les mêmes résultats. Les compteurs d'exécution trahissent l'\omterm{def:pl:the-research-environment:hidden}{état caché} --- un compteur sauté signifie une
cellule relancée ou exécutée puis supprimée, un compteur inférieur à celui de la cellule qui la précède signifie une exécution hors de l'ordre du document --- mais seulement si quelqu'un
les lit.

Le notebook piégé du chapitre calcule le chiffre phare du débat d'ouverture sur un panel synthétique de vingt pays entre 1946 et 2009~: la
moyenne de la croissance réelle des observations pays-année dont la dette dépasse 90~\% du PIB, moyennée par pays. Son historique est consigné dans le code~: les
cellules, dans l'ordre du document, et ce qui a été exécuté, dans l'ordre où il l'a été.

\omcode{code/platforms/15-the-research-environment/python/pl_envaudit.py}{44}{53}{Le notebook piégé~: cinq cellules de code dans l'ordre du document,
une cellule qui a été exécutée puis supprimée, et l'ordre dans lequel l'analyste les a exécutées --- la cellule de nettoyage, proche du haut, en dernier.}\label{lst:pl:the-research-environment:planted}

\begin{omfigure}
\begin{tikzpicture}[font=\scriptsize,cell/.style={draw,rounded corners=2pt,minimum width=5.6cm,minimum height=0.55cm,anchor=west,
  fill=omMeth!10,font=\ttfamily\scriptsize},>=Stealth]
\foreach \y/\n/\s in {0/1/{df = pd.read\_csv(\textquotedbl{}panel.csv\textquotedbl{})},-0.8/6/{df = df[df.year >= 1950]},-1.6/2/{high = df[df.debt > 90]},
  -2.4/4/{sel = high[\textasciitilde high.country.isin(exclude)]},-3.2/5/{result = sel.groupby(\dots)\dots{} -> 1.96}} {
  \node[cell] at (0,\y) {\s};
  \node[anchor=east,font=\scriptsize] at (-0.1,\y) {[\n]};}
\node[cell,fill=omThm!12,dashed] (del) at (7.2,-2.0) {exclude = [\textquotedbl{}C01\textquotedbl{}, \dots, \textquotedbl{}C05\textquotedbl{}]};
\node[anchor=west,font=\scriptsize\itshape,omThm] at (7.2,-2.6) {exécutée en [3], puis supprimée};
\node[font=\scriptsize\itshape] at (2.8,0.65) {ordre du document, avec les compteurs d'exécution};
\draw[->,omThm,thick] (5.7,-0.8) .. controls (6.6,-1.5) and (6.6,-3.0) .. (5.7,-3.2);
\node[omThm,font=\scriptsize,anchor=west] at (6.5,-3.4) {exécutée après le résultat};
\end{tikzpicture}

\omcaption{Le notebook piégé tel qu'enregistré~: cellules de code dans l'ordre du document, avec leurs compteurs d'exécution. La cellule de nettoyage, deuxième de la
page, a été exécutée en sixième~; la cellule qui a défini \texttt{exclude} a été exécutée en troisième puis supprimée. Le 1.96 enregistré provient de cet
historique.}\label{fig:pl:the-research-environment:notebook}
\end{omfigure}

Enregistré, le notebook affiche 1.96 sous sa cellule de résultat. Rien dans le document n'indique que cinq pays ont été exclus, ni que les années antérieures
à 1950 figuraient encore dans les données au moment où le chiffre a été calculé.

\section{Audit du document}

\texttt{firm.envaudit} lit le JSON du notebook et rien d'autre (\cref{lst:pl:the-research-environment:audit}). Il parcourt les cellules de code
dans l'ordre du document avec leurs compteurs, et analyse chaque cellule pour voir quels noms elle lit et lesquels elle lie.

\omcode{code/firm/envaudit/firm_envaudit.py}{89}{118}{L'audit~: compteurs hors d'ordre du document, compteurs sautés, sorties sans exécution,
cellules jamais exécutées, et noms lus avant qu'une cellule située au-dessus ne les lie --- ou liés par aucune cellule.}\label{lst:pl:the-research-environment:audit}

Sur le notebook piégé, il relève trois anomalies (\cref{tab:pl:the-research-environment:audit}). La cellule qui sélectionne les observations pays-année au-dessus de 90~\%
a été exécutée en deuxième, mais une cellule située au-dessus d'elle a été exécutée en sixième~: la cellule de nettoyage a été lancée après le résultat, si bien que les données affichées ne sont pas celles que le résultat
a utilisées. Le compteur~3 manque~: une cellule a été exécutée, puis elle a disparu. Et la cellule de sélection lit \texttt{exclude}, que nulle cellule du notebook ne définit.

\begin{table}[ht]
\centering\small
\begin{tabular}{llp{7.4cm}}
\toprule
constat & cellule & détail \\
\midrule
hors d'ordre & 3 & compteur 2 inférieur à celui de la cellule au-dessus (6, cellule 2) \\
compteur sauté & --- & compteur 3 manquant~: une cellule relancée ou supprimée \\
jamais défini & 4 & \texttt{exclude} \\
\bottomrule
\end{tabular}
\caption{L'audit du notebook piégé, à partir de son seul JSON (la cellule 0 est le titre).}\label{tab:pl:the-research-environment:audit}
\end{table}

Relancé depuis le début dans un interpréteur neuf (\cref{lst:pl:the-research-environment:rerun}), le notebook échoue~:
\texttt{NameError: name 'exclude' is not defined}. Le chiffre affiché ne peut pas être produit par la page elle-même.

\omcode{code/firm/envaudit/firm_envaudit.py}{132}{140}{La relance~: les cellules de code, du haut vers le bas, en un seul script, dans un interpréteur neuf lancé dans un
sous-processus.}\label{lst:pl:the-research-environment:rerun}

\section{Du notebook à la bibliothèque}

\begin{definition}[Bibliothèque de recherche]\label{def:pl:the-research-environment:library}
Une \emph{bibliothèque de recherche}\index{bibliothèque de recherche} est le code partagé, versionné et testé qu'importent les notebooks et les chaînes de traitement de recherche~:
des calculs extraits des notebooks dès qu'ils servent plus d'une fois, dont les choix sont des arguments explicites et dont le comportement
est fixé par des tests.
\end{definition}

La solution consiste à décider, au grand jour, ce que le calcul est. Sans exclusion et avec les cellules exécutées telles qu'elles sont, le notebook affiche 2.44. Le
calcul appartient à la bibliothèque (\cref{lst:pl:the-research-environment:library})~: le seuil, la première année, les pays exclus et
la pondération deviennent des arguments, et chacun fait l'objet d'un test. La bibliothèque reproduit les trois chiffres sur lesquels porte le débat~: 1.96 avec les cinq pays
exclus et toutes les années (le chiffre enregistré, expliqué), 2.44 avec tous les pays à partir de 1950 (le chiffre corrigé), et 2.90 si chaque
observation pays-année compte une fois au lieu de chaque pays --- le choix de pondération que la critique de l'article d'ouverture a également mis en cause. Les cinq pays exclus
représentent 204 des 455 observations pays-année au-dessus de 90~\% depuis 1950.

\begin{omfigure}
\begin{tikzpicture}
\begin{axis}[width=12cm,height=5.6cm,xbar,bar width=7pt,xlabel={croissance moyenne au-dessus de 90~\% de dette (\%)},xmin=0,xmax=3.6,
  symbolic y coords={all countries; from 1950,all countries; all years,five excluded; from 1950,five excluded; all years},yticklabels={cinq pays exclus~; toutes les années,cinq pays exclus~; depuis 1950,tous les pays~; toutes les années,tous les pays~; depuis 1950},
  ytick=data,enlarge y limits=0.2,tick label style={font=\small},label style={font=\small},grid=major,grid style={gray!20},
  table/col sep=comma,nodes near coords,nodes near coords style={font=\scriptsize,/pgf/number format/fixed,/pgf/number format/fixed zerofill,/pgf/number format/precision=2},
  legend style={font=\footnotesize,at={(0.5,1.03)},anchor=south,legend columns=2}]
\addplot[fill=omDef!45,draw=omDef] table[y=case,x=by_country]{figdata/platforms/15-the-research-environment/choices.csv};
\addlegendentry{moyenne des moyennes par pays}
\addplot[fill=omMeth!45,draw=omMeth] table[y=case,x=by_country_year]{figdata/platforms/15-the-research-environment/choices.csv};
\addlegendentry{moyenne des observations pays-année}
\end{axis}
\end{tikzpicture}

\omcaption{Le résultat principal sous les trois choix que le notebook dissimulait~: exclure ou non cinq pays, partir de 1946 ou de 1950, et moyenner les
pays ou les observations pays-année. Le 1.96 enregistré est la moyenne des moyennes par pays de la paire du haut (en bleu)~; le 2.44 corrigé, celle de la paire du bas. Données~:
\texttt{fig\_envaudit.py}.}\label{fig:pl:the-research-environment:choices}
\end{omfigure}

\omcode{code/platforms/15-the-research-environment/python/research_lib.py}{8}{24}{La fonction promue~: les choix qui vivaient dans la
mémoire du notebook sont des arguments.}\label{lst:pl:the-research-environment:library}

\section{Revue de code}

\begin{definition}[Revue de code]\label{def:pl:the-research-environment:review}
La \emph{revue de code}\index{revue de code} est l'examen d'une modification apportée à du code partagé par une personne autre que son auteur, avant que celle-ci soit fusionnée~:
de ce que fait la modification, si ses tests le montrent et si elle fait bien ce que dit sa description.
\end{definition}

La \omterm{def:pl:the-research-environment:review}{revue de code} est le lieu où la revue de recherche du livre~7 (chapitre~1) rencontre la base de code. Un relecteur ne peut pas examiner l'\omterm{def:pl:the-research-environment:hidden}{état caché} d'un notebook, seulement un
diff~: c'est la promotion d'un calcul vers la bibliothèque qui le rend examinable. La liste de contrôle de promotion du chapitre, exécutée par
\texttt{promotion\_checklist}, est le minimum qu'exige un relecteur~: l'audit ne trouve aucun \omterm{def:pl:the-research-environment:hidden}{état caché}, le notebook se réexécute depuis le haut, la réexécution
reproduit le résultat enregistré, les fonctions de la bibliothèque sont testées, l'environnement correspond au verrouillage de l'environnement, les données correspondent à leur instantané de données.

\begin{definition}[Monorepo]\label{def:pl:the-research-environment:monorepo}
Un \emph{monorepo}\index{monorepo} est un dépôt unique de gestion de versions qui regroupe le code de nombreux projets, bibliothèques et équipes,
de sorte qu'une modification d'une bibliothèque et de chacun de ses appelants puisse être faite, relue et testée ensemble.
\end{definition}

Le retour d'expérience de Google (Potvin et Levenberg, 2016) décrit un unique dépôt servant de source commune de vérité pour des dizaines de
milliers de développeurs, rendu praticable par des systèmes et des flux de travail conçus pour sa taille. Le code propre à cette série (\texttt{code/firm/}) est un \omterm{def:pl:the-research-environment:monorepo}{monorepo} en miniature~: un
composant d'un livre est utilisé par le suivant, et toute modification de celui-ci est testée sur chaque chapitre qui l'importe.

\section{Environnements et conteneurs}

\begin{definition}[Image de conteneur]\label{def:pl:the-research-environment:container}
Une \emph{image de conteneur}\index{image de conteneur} est un système de fichiers et une configuration empaquetés --- une couche de système d'exploitation, l'interpréteur,
les bibliothèques, le code --- à partir desquels des environnements isolés et identiques peuvent être lancés sur toute machine qui exécute des conteneurs.
\end{definition}

L'Open Container Initiative précise le format~: un manifeste d'image, un index d'image facultatif, un ensemble de couches de système de fichiers et une
configuration, de sorte que les outils de construction, de transport et d'exécution des images soient interopérables. Un conteneur fige un environnement~; un
verrouillage de l'environnement (livre~7, chapitre~29) en décrit un. Le chapitre utilise le plus léger des deux~: \texttt{capture\_env} énumère l'interpréteur et chaque
distribution installée, et \texttt{compare\_env} les confronte au fichier \texttt{requirements.txt} du dépôt~; les données sont vérifiées
à l'aide d'un manifeste d'empreintes (\texttt{snapshot}). Là où ce verrouillage suffit, un conteneur n'ajoute que du poids~; là où les bibliothèques système ou
les compilateurs comptent (les \omterm{def:pl:native-extensions:extension}{extensions natives} du chapitre~9), c'est la seule trace honnête.

\begin{dated}{2026-09}{Outils de reproductibilité}\label{dat:pl:the-research-environment:tools}
Le format de notebook Jupyter (nbformat) enregistre le \texttt{execution\_count} de chaque cellule de code sous forme d'entier ou de valeur nulle, et ses sorties sous forme de liste~;
la spécification d'image de l'Open Container Initiative définit une image comme un manifeste, un index facultatif, des couches de système de fichiers et une
configuration. Ce sont deux spécifications ouvertes~; les linters de notebooks et les outils de construction de conteneurs qui s'appuient dessus changent souvent, les formats beaucoup moins.
\end{dated}

\section{La reproductibilité en routine}

La dernière étape reproduit le chiffre corrigé à partir d'un répertoire propre~: la bibliothèque promue, le fichier de données et un script de trois lignes qui
importe la bibliothèque et affiche le résultat principal. Le fichier de données est contrôlé par rapport au hachage de l'instantané (aucune différence), l'environnement
confronté au verrouillage (aucune différence), et le script affiche 2.44 dans un interpréteur neuf. C'est le résultat reproductible du livre~7 (chapitre~29), obtenu
non par soin mais par routine~: un résultat n'est publié qu'avec la commande qui le recalcule à partir d'une copie de travail propre.

\begin{omfigure}
\begin{tikzpicture}[font=\scriptsize,box/.style={draw,rounded corners=2pt,align=center,minimum height=0.9cm,minimum width=2.2cm,
  fill=omDef!10},>=Stealth]
\node[box] (a) at (0,0) {répertoire\\ propre};
\node[box] (b) at (3.0,0) {fichier de données\\ contre hachage\\ de l'instantané};
\node[box] (c) at (6.0,0) {environnement\\ contre verrouillage};
\node[box] (d) at (9.0,0) {interpréteur neuf~:\\ bibliothèque + script};
\node[box,fill=omMeth!15,minimum width=1.2cm] (e) at (11.4,0) {2.44};
\foreach \x/\y in {a/b,b/c,c/d,d/e} {\draw[->] (\x) -- (\y);}
\node[font=\scriptsize\itshape] at (3.0,-0.95) {aucune différence};
\node[font=\scriptsize\itshape] at (6.0,-0.95) {aucune différence};
\end{tikzpicture}

\omcaption{La reproduction en routine~: un répertoire propre avec la bibliothèque promue et le fichier de données, les données vérifiées par rapport au
hachage de l'instantané et l'environnement au verrouillage, et le résultat principal recalculé dans un interpréteur neuf.}\label{fig:pl:the-research-environment:repro}
\end{omfigure}

\begin{example}[Réexécution depuis le haut]\label{ex:pl:the-research-environment:rerun}
Le notebook piégé affiche 1.96. L'audit repère une cellule exécutée hors d'ordre, une cellule supprimée et un nom défini nulle part~; la réexécution depuis le haut
échoue sur ce nom. Sans l'exclusion, et avec les cellules telles qu'elles sont, le résultat est 2.44~; la bibliothèque montre que 1.96 est exactement la
moyenne obtenue en excluant cinq pays et en gardant les années antérieures à 1950 dans les données, et que regrouper les observations pays-année au lieu des pays donnerait
2.90. Depuis un répertoire propre, avec l'instantané de données et le verrouillage de l'environnement vérifiés, la bibliothèque reproduit 2.44.
\end{example}

\section{Tutoriel~: audit, réexécution, promotion, reproduction}\label{tut:pl:the-research-environment:tut}

\begin{tutorial}
\textbf{Objectif.} Repérer l'\omterm{def:pl:the-research-environment:hidden}{état caché} d'un notebook à partir de son JSON, montrer que son résultat ne peut pas être recalculé, corriger et promouvoir le calcul, et
le reproduire depuis un répertoire propre. \textbf{État final~:} le \cref{tab:pl:the-research-environment:audit} et les trois contrôles de la reproduction.
\begin{tutsteps}
\item \textbf{Piégeage}~: \texttt{pl\_envaudit.planted(dir)} écrit les données et le notebook avec les sorties de son historique.
\item \textbf{Audit}~: \texttt{firm\_envaudit.audit(nb)}.
\item \textbf{Réexécution}~: \texttt{firm\_envaudit.rerun(nb, dir)}, puis avec \texttt{fixed(nb)}.
\item \textbf{Promotion}~: \texttt{research\_lib.average\_growth} et ses tests.
\item \textbf{Reproduction}~: \texttt{pl\_envaudit.reproduce(clean\_dir, dir)}~: instantané, verrouillage, réexécution.
\end{tutsteps}
\textbf{Pour aller plus loin.} Appliquer l'audit à ses propres notebooks~; ajouter un contrôle qui vérifie que chaque sortie enregistrée d'un notebook est reproduite par sa
réexécution (sorties comparées cellule par cellule).
\end{tutorial}

\section{Construction~: l'audit de l'environnement}\label{bld:pl:the-research-environment:envaudit}

\begin{build}
\bfield{Rôle} Le verrou de la plateforme entre l'exploration et tout ce sur quoi d'autres s'appuient~: le chiffre d'un notebook ne devient un résultat que lorsqu'il peut
être recalculé depuis le haut, à partir d'une copie de travail propre, dans l'environnement verrouillé, sur l'instantané de ses données.
\bfield{Interface} \texttt{load}, \texttt{code\_cells}, \texttt{audit -> [Finding]}, \texttt{saved\_output}, \texttt{rerun -> Rerun},
\texttt{capture\_env}, \texttt{compare\_env}, \texttt{snapshot}, \texttt{check\_snapshot}, \texttt{promotion\_checklist}.
\bfield{Règles} L'audit ne lit que le document~; la réexécution utilise un interpréteur neuf dans un sous-processus, jamais le noyau de l'auteur~; les environnements
sont comparés au fichier de verrouillage, les données aux empreintes.
\bfield{Tests d'acceptation} \texttt{code/firm/envaudit/tests/}~: un notebook propre ne produit aucun constat et se réexécute~; un notebook piégé produit chaque type de
constat et une réexécution échouée~; les imports, les fonctions et les fonctions intégrées ne sont pas des constats~; les comparaisons au verrouillage et à l'instantané~; la liste de contrôle réussit et échoue
comme il se doit.
\bfield{Pour aller plus loin} Comparaison cellule par cellule des sorties enregistrées et réexécutées~; une construction de conteneur à partir du verrouillage~; l'audit en tant que hook de pre-commit dans le
\omterm{def:pl:the-research-environment:monorepo}{monorepo}.
\end{build}

\begin{omsources}
\item T.~Herndon, M.~Ash et R.~Pollin, ``Does high public debt consistently stifle economic growth? A critique of Reinhart and Rogoff'',
\emph{Cambridge Journal of Economics} 38(2), 2014 (PERI working paper 322, 2013).
\item J.~F. Pimentel, L.~Murta, V.~Braganholo et J.~Freire, ``A large-scale study about quality and reproducibility of Jupyter
notebooks'', MSR 2019.
\item R.~Potvin et J.~Levenberg, ``Why Google stores billions of lines of code in a single repository'', \emph{Communications of the ACM}
59(7), 2016.
\item Jupyter nbformat documentation~; Open Container Initiative image specification.
\end{omsources}

\section{Exercices}

\begin{exercise}[$\star$]\label{exo:pl:the-research-environment:1}
Un notebook affiche, dans l'ordre du document, les compteurs 1, 2, 5, 3, 6 de ses cellules de code. Que signale l'audit~?
\end{exercise}

\begin{exercise}[$\star$]\label{exo:pl:the-research-environment:2}
Pourquoi l'audit signale-t-il \texttt{exclude} comme jamais défini plutôt que comme utilisé avant d'être défini~?
\end{exercise}

\begin{exercise}[$\star$]\label{exo:pl:the-research-environment:3}
Quelles sont les deux différences entre le calcul enregistré et le calcul corrigé qui expliquent l'écart entre 1.96 et 2.44~?
\end{exercise}

\begin{exercise}[$\star\star$]\label{exo:pl:the-research-environment:4}
Un notebook dont les compteurs se suivent dans l'ordre du document peut-il encore avoir un \omterm{def:pl:the-research-environment:hidden}{état caché}~? Donner un exemple.
\end{exercise}

\begin{exercise}[$\star\star$]\label{exo:pl:the-research-environment:5}
Pourquoi pondérer par pays (2.44) ou par pays-année (2.90) est-il un choix qui relève des arguments de la fonction, et lequel retenir pour le rapport~?
\end{exercise}

\begin{exercise}[$\star\star$]\label{exo:pl:the-research-environment:6}
Quand un verrouillage de l'environnement suffit-il, et quand faut-il une \omterm{def:pl:the-research-environment:container}{image de conteneur}~?
\end{exercise}

\begin{exercise}[$\star\star\star$]\label{exo:pl:the-research-environment:7}
\emph{Programmation.} Ajouter à \texttt{firm.envaudit} un contrôle qui compare la sortie enregistrée de chaque cellule à la sortie de la réexécution pour cette cellule. Que
signale-t-il sur le notebook corrigé, et pourquoi est-il plus difficile que la comparaison de la dernière sortie~?
\end{exercise}

\begin{exercise}[$\star\star\star$]\label{exo:pl:the-research-environment:8}
\emph{Trouver la faille.} «~Le notebook est versionné et le relecteur l'a approuvé, donc son résultat est reproductible.~»
\end{exercise}

\section{Problème~: réexécuter depuis le haut}

\begin{problem}[{Problème du week-end --- un nombre que la page ne peut pas produire}]\label{pb:pl:the-research-environment:1}
Le notebook piégé du \cref{lst:pl:the-research-environment:planted}, \texttt{firm.envaudit} et la bibliothèque promue.

\medskip\noindent\textbf{Partie I --- Le notebook.}
\begin{enumerate}
\item Que stocke un notebook Jupyter pour chaque cellule de code~?
\item Dans quel ordre les cellules ont-elles été exécutées, et dans quel ordre sont-elles présentées~?
\item Quel nombre le notebook affiche-t-il~?
\item Qu'est-ce que l'\omterm{def:pl:the-research-environment:hidden}{état caché}, sous trois formes~?
\item Que rapporte l'étude à grande échelle des notebooks sur GitHub~?
\end{enumerate}

\medskip\noindent\textbf{Partie II --- L'audit.}
\begin{enumerate}[resume]
\item Quels trois constats l'audit établit-il, et à partir de quoi~?
\item Que fait la réexécution depuis le haut~?
\item Qu'est-ce que chaque constat explique du nombre enregistré~?
\item Que ne voit pas l'audit~?
\item Pourquoi réexécuter dans un interpréteur neuf, dans un sous-processus~?
\end{enumerate}

\medskip\noindent\textbf{Partie III --- La bibliothèque.}
\begin{enumerate}[resume]
\item Qu'est-ce qui devient un argument de la fonction promue~?
\item Quels sont les trois nombres qu'elle reproduit, et selon quels choix~?
\item Combien de pays-années au-dessus de 90~\% depuis 1950 détiennent les cinq pays exclus~?
\item Que vérifie un relecteur avant que la promotion ne soit fusionnée~?
\item Pourquoi un \omterm{def:pl:the-research-environment:monorepo}{monorepo} est-il utile~?
\end{enumerate}

\medskip\noindent\textbf{Partie IV --- Le verdict.}
\begin{enumerate}[resume]
\item Énoncer le \emph{résultat nommé}~: le nombre phare tel qu'enregistré et tel que réexécuté, les constats qui expliquent l'écart, et la reproduction
du résultat corrigé depuis un environnement propre.
\item Que faudrait-il exiger avant qu'un nombre issu d'un notebook entre dans un rapport de recherche~?
\item Quand faut-il construire une \omterm{def:pl:the-research-environment:container}{image de conteneur}~?
\item Que devrait consigner un journal de recherche (livre~7, chapitre~1) sur cette analyse~?
\item En une phrase~: qu'est-ce qui rend un nombre reproductible~?
\end{enumerate}
\end{problem}

\section{Questions d'entretien}

\begin{interviewq}[$\star$ \iqroles{chercheur}]\label{iq:pl:the-research-environment:1}
Le notebook d'un collègue affiche un résultat que vous ne parvenez pas à reproduire. Que vérifiez-vous~?
\end{interviewq}

\begin{interviewq}[$\star\star$ \iqroles{chercheur, développeur}]\label{iq:pl:the-research-environment:2}
Quand le code doit-il passer d'un notebook à une bibliothèque~?
\end{interviewq}

\begin{interviewq}[$\star\star$ \iqroles{développeur}]\label{iq:pl:the-research-environment:3}
Qu'est-ce qu'un conteneur vous apporte que ne fait pas un fichier de dépendances, et que coûte-t-il~?
\end{interviewq}

\begin{interviewq}[$\star\star$ \iqroles{chercheur}]\label{iq:pl:the-research-environment:4}
Comment rendriez-vous reproductible, un an plus tard, chaque nombre d'un rapport de recherche~?
\end{interviewq}

\begin{interviewq}[$\star\star\star$ \iqroles{développeur, chercheur}]\label{iq:pl:the-research-environment:5}
Concevez l'environnement de recherche d'une équipe de vingt chercheurs qui partagent données et code.
\end{interviewq}
